\documentclass[12pt, a4paper]{article}

% ── Packages ─────────────────────────────────────────────────────────────────
\setlength{\headheight}{15pt}
\usepackage[french]{babel}
\usepackage[utf8]{inputenc}
\usepackage[T1]{fontenc}
\usepackage[margin=2.5cm, top=3cm]{geometry}
\usepackage{hyperref}
\usepackage{xcolor}
\usepackage{amssymb}
\usepackage{listings}
\usepackage{titlesec}
\usepackage{enumitem}
\usepackage{fancyhdr}
\usepackage{graphicx}
\usepackage{tcolorbox}
\usepackage{fontawesome5}
\usepackage{booktabs}
\usepackage{amsmath}

% ── Colors ───────────────────────────────────────────────────────────────────
\definecolor{teal}{RGB}{20, 184, 166}
\definecolor{indigo}{RGB}{99, 102, 241}
\definecolor{codebg}{RGB}{30, 30, 30}
\definecolor{codetext}{RGB}{220, 220, 220}
\definecolor{infobg}{RGB}{240, 249, 255}
\definecolor{infoborder}{RGB}{14, 165, 233}
\definecolor{warnbg}{RGB}{255, 251, 235}
\definecolor{warnborder}{RGB}{245, 158, 11}
\definecolor{successbg}{RGB}{240, 253, 244}
\definecolor{successborder}{RGB}{34, 197, 94}

% ── Hyperref config ───────────────────────────────────────────────────────────
\hypersetup{
  colorlinks=true,
  linkcolor=teal,
  urlcolor=indigo,
  pdftitle={Guide de Déploiement MediCloud},
  pdfauthor={Noaman Makhlouf}
}

% ── Code listings style ───────────────────────────────────────────────────────
\lstset{
  backgroundcolor=\color{codebg},
  basicstyle=\ttfamily\small\color{codetext},
  breaklines=true,
  frame=single,
  framesep=5pt,
  rulecolor=\color{teal},
  xleftmargin=10pt,
  xrightmargin=10pt,
  aboveskip=10pt,
  belowskip=10pt,
  showstringspaces=false,
  keywordstyle=\color{teal}\bfseries,
  commentstyle=\color{gray}\itshape,
  stringstyle=\color{orange},
}

% ── Custom boxes ─────────────────────────────────────────────────────────────
\tcbuselibrary{breakable, skins}

\newtcolorbox{infobox}[1]{
  colback=infobg,
  colframe=infoborder,
  fonttitle=\bfseries,
  title={#1},
  breakable
}

\newtcolorbox{warnbox}[1]{
  colback=warnbg,
  colframe=warnborder,
  fonttitle=\bfseries,
  title={#1},
  breakable
}

\newtcolorbox{successbox}[1]{
  colback=successbg,
  colframe=successborder,
  fonttitle=\bfseries,
  title={#1},
  breakable
}

% ── Section styles ────────────────────────────────────────────────────────────
\titleformat{\section}
  {\large\bfseries\color{teal}}
  {\thesection.}{0.5em}{}[\titlerule]

\titleformat{\subsection}
  {\normalsize\bfseries\color{indigo}}
  {\thesubsection.}{0.5em}{}

% ── Header/Footer ─────────────────────────────────────────────────────────────
\pagestyle{fancy}
\fancyhf{}
\rhead{\textcolor{teal}{\textbf{MediCloud}} — Guide de Déploiement}
\lhead{\textcolor{gray}{\small Noaman Makhlouf}}
\rfoot{\thepage}
\lfoot{\textcolor{gray}{\small ENSEM — Application Cloud Native}}

% ─────────────────────────────────────────────────────────────────────────────
\begin{document}

% ── Title page ───────────────────────────────────────────────────────────────
\begin{titlepage}
  \centering
  \vspace*{2cm}

  {\huge\bfseries\color{teal} MediCloud}\\[0.5cm]
  {\Large\color{indigo} Guide Complet de Déploiement}\\[0.3cm]
  {\large\color{gray} Du clonage GitHub au déploiement Vercel}

  \vspace{1.5cm}
  \rule{0.6\textwidth}{1.5pt}
  \vspace{1cm}

  \begin{tcolorbox}[colback=teal!10, colframe=teal, width=0.8\textwidth]
    \centering
    Ce guide a été rédigé pour permettre à n'importe quel développeur de
    \textbf{lancer MediCloud localement} et de le \textbf{déployer en production}
    sur Vercel, pas à pas, depuis zéro.
  \end{tcolorbox}

  \vspace{1.5cm}
  {\textbf{Auteur :} Noaman Makhlouf}\\[0.3cm]
  {\textbf{Dépôt GitHub :} \href{https://github.com/Noamane2023/medicloud}{\texttt{github.com/Noamane2023/medicloud}}}\\[0.3cm]
  {\textbf{Stack technique :} Next.js 16 · Supabase · Upstash Redis · Vercel}\\[0.3cm]
  {\textbf{Date :} \today}
\end{titlepage}

\tableofcontents
\newpage

% ─────────────────────────────────────────────────────────────────────────────
\section{Introduction — Qu'est-ce que MediCloud ?}
% ─────────────────────────────────────────────────────────────────────────────

\textbf{MediCloud} est une plateforme de télémédecine Cloud-Native développée
dans le cadre d'un projet académique à l'ENSEM. Elle permet à des patients
de trouver un médecin, de réserver un rendez-vous en ligne, et de gérer
un dossier médical numérique — le tout depuis un smartphone ou un navigateur.

\subsection*{Architecture technique en bref}

\begin{center}
\begin{tabular}{lll}
  \toprule
  \textbf{Couche} & \textbf{Technologie} & \textbf{Rôle} \\
  \midrule
  Front-end  & Next.js 16 + React & Interface utilisateur PWA \\
  Back-end   & API Routes (Serverless) & Logique métier \\
  Base de données & Supabase (PostgreSQL) & Données persistantes \\
  Cache      & Upstash Redis & Réponses rapides sous charge \\
  Emails     & Resend API & Rappels de rendez-vous \\
  Déploiement & Vercel & Hébergement Cloud Edge \\
  \bottomrule
\end{tabular}
\end{center}

% ─────────────────────────────────────────────────────────────────────────────
\section{Étape 1 — Prérequis : Logiciels à Installer}
% ─────────────────────────────────────────────────────────────────────────────

Avant de commencer, tu dois avoir installé les outils suivants sur ton
ordinateur (Windows, macOS ou Linux).

\subsection{Node.js (v18 ou supérieur)}

Node.js est le moteur d'exécution JavaScript utilisé par Next.js.

\begin{enumerate}[label=\arabic*.]
  \item Rends-toi sur \href{https://nodejs.org}{\texttt{nodejs.org}}
  \item Télécharge la version \textbf{LTS} (Long Term Support)
  \item Installe-le en gardant les options par défaut
  \item Vérifie l'installation en ouvrant un terminal et en tapant :
\end{enumerate}

\begin{lstlisting}[language=bash]
node --version
# Attendu : v18.x.x ou superieur

npm --version
# Attendu : 9.x.x ou superieur
\end{lstlisting}

\subsection{Git}

Git est le système de contrôle de version utilisé pour télécharger le code
depuis GitHub.

\begin{enumerate}[label=\arabic*.]
  \item Rends-toi sur \href{https://git-scm.com/downloads}{\texttt{git-scm.com/downloads}}
  \item Télécharge et installe Git selon ton système d'exploitation
  \item Vérifie l'installation :
\end{enumerate}

\begin{lstlisting}[language=bash]
git --version
# Attendu : git version 2.x.x
\end{lstlisting}

\subsection{Visual Studio Code (éditeur recommandé)}

\href{https://code.visualstudio.com}{\texttt{code.visualstudio.com}} — Télécharge
et installe VS Code. C'est l'éditeur de code le plus populaire pour le
développement web moderne.

\textbf{Extensions recommandées à installer dans VS Code :}
\begin{itemize}
  \item \texttt{ESLint} — Détection d'erreurs JavaScript/TypeScript
  \item \texttt{Tailwind CSS IntelliSense} — Autocomplétion CSS
  \item \texttt{Prettier} — Formatage automatique du code
  \item \texttt{GitLens} — Visualisation Git avancée
\end{itemize}

% ─────────────────────────────────────────────────────────────────────────────
\section{Étape 2 — Créer un Compte GitHub et Accéder au Dépôt}
% ─────────────────────────────────────────────────────────────────────────────

\subsection{Créer un compte GitHub}

\begin{enumerate}[label=\arabic*.]
  \item Va sur \href{https://github.com}{\texttt{github.com}}
  \item Clique sur \textbf{"Sign up"}
  \item Remplis le formulaire : adresse email, mot de passe, nom d'utilisateur
  \item Confirme ton email
  \item Envoie ton \textbf{nom d'utilisateur GitHub} à Noaman pour être ajouté
        comme collaborateur sur le dépôt
\end{enumerate}

\begin{warnbox}{Important — Communique ton pseudonyme GitHub}
  Dis à Noaman ton nom d'utilisateur GitHub (pas ton email) pour qu'il
  t'ajoute aux collaborateurs du dépôt. Tu recevras une invitation par email
  à accepter.
\end{warnbox}

\subsection{Accepter l'invitation de collaboration}

\begin{enumerate}[label=\arabic*.]
  \item Vérifie ta boîte email
  \item Clique sur le lien d'\textbf{invitation GitHub}
  \item Accepte l'accès au dépôt \texttt{Noamane2023/medicloud}
\end{enumerate}

% ─────────────────────────────────────────────────────────────────────────────
\section{Étape 3 — Cloner et Lancer le Projet Localement}
% ─────────────────────────────────────────────────────────────────────────────

\subsection{Cloner le dépôt}

Ouvre un terminal (PowerShell, Terminal macOS ou Bash Linux) et exécute :

\begin{lstlisting}[language=bash]
# Navigue vers le dossier de ton choix
cd Documents

# Clone le depot depuis GitHub
git clone https://github.com/Noamane2023/medicloud.git

# Entre dans le dossier du projet
cd medicloud
\end{lstlisting}

\subsection{Installer les dépendances}

\begin{lstlisting}[language=bash]
# Installe tous les packages necessaires
npm install
\end{lstlisting}

Cette commande va télécharger toutes les bibliothèques listées dans le
fichier \texttt{package.json} (React, Next.js, Supabase SDK, etc.).

\subsection{Configurer les variables d'environnement}

Le projet a besoin de clés secrètes pour se connecter aux services en ligne.
Ces clés ne sont \textbf{jamais} publiées sur GitHub pour des raisons de sécurité.

\begin{enumerate}[label=\arabic*.]
  \item Crée un fichier \texttt{.env.local} à la racine du projet
  \item Demande à Noaman de te partager les valeurs privées
  \item Remplis le fichier ainsi :
\end{enumerate}

\begin{lstlisting}
# .env.local - Ne jamais partager ce fichier !

NEXT_PUBLIC_SUPABASE_URL=https://xxxxx.supabase.co
NEXT_PUBLIC_SUPABASE_ANON_KEY=eyJhbGciOiJ...
SUPABASE_SERVICE_ROLE_KEY=eyJhbGciOiJ...

RESEND_API_KEY=re_xxxxxxxxxxxx
CRON_SECRET=medicloud_cron_2026

UPSTASH_REDIS_REST_URL=https://xxx.upstash.io
UPSTASH_REDIS_REST_TOKEN=xxxxxxxx
\end{lstlisting}

\subsection{Lancer le serveur de développement}

\begin{lstlisting}[language=bash]
npm run dev
\end{lstlisting}

\begin{successbox}{Succès !}
  Ouvre ton navigateur sur \textbf{\href{http://localhost:3000}{\texttt{http://localhost:3000}}}.
  Tu devrais voir la page d'accueil de MediCloud. \\
  Le serveur local se rechargera automatiquement à chaque modification du code.
\end{successbox}

% ─────────────────────────────────────────────────────────────────────────────
\section{Étape 4 — Déploiement en Production sur Vercel}
% ─────────────────────────────────────────────────────────────────────────────

Vercel est la plateforme Cloud officielle de déploiement pour les applications
Next.js. Elle offre un \textbf{déploiement continu} : chaque push sur GitHub
met à jour automatiquement l'application en production.

\subsection{Créer un compte Vercel}

\begin{enumerate}[label=\arabic*.]
  \item Va sur \href{https://vercel.com}{\texttt{vercel.com}}
  \item Clique sur \textbf{"Sign Up"}
  \item Authentifie-toi avec ton \textbf{compte GitHub} (recommandé)
\end{enumerate}

\subsection{Importer le projet depuis GitHub}

\begin{enumerate}[label=\arabic*.]
  \item Dans le Dashboard Vercel, clique sur \textbf{"Add New... → Project"}
  \item Sélectionne le dépôt \texttt{medicloud} depuis la liste GitHub
  \item Vercel détecte automatiquement qu'il s'agit d'un projet \textbf{Next.js}
  \item \textbf{Ne clique pas encore sur "Deploy"} — il faut d'abord ajouter les variables
\end{enumerate}

\subsection{Ajouter les variables d'environnement}

Dans la section \textbf{"Environment Variables"} de la page de configuration Vercel :

\begin{center}
\begin{tabular}{ll}
  \toprule
  \textbf{Nom de la variable} & \textbf{Description} \\
  \midrule
  \texttt{NEXT\_PUBLIC\_SUPABASE\_URL} & URL de ta base de données Supabase \\
  \texttt{NEXT\_PUBLIC\_SUPABASE\_ANON\_KEY} & Clé publique Supabase \\
  \texttt{SUPABASE\_SERVICE\_ROLE\_KEY} & Clé admin secrète Supabase \\
  \texttt{RESEND\_API\_KEY} & Clé API pour l'envoi d'emails \\
  \texttt{CRON\_SECRET} & Mot de passe pour sécuriser les Cron Jobs \\
  \texttt{UPSTASH\_REDIS\_REST\_URL} & URL de la base de données Redis \\
  \texttt{UPSTASH\_REDIS\_REST\_TOKEN} & Token d'authentification Redis \\
  \bottomrule
\end{tabular}
\end{center}

\begin{warnbox}{Attention à la sécurité}
  Les variables préfixées \texttt{NEXT\_PUBLIC\_} sont visibles par le navigateur.
  Ne mets \textbf{jamais} de données sensibles (clés secrètes) dans des variables
  \texttt{NEXT\_PUBLIC\_}. Supabase et Vercel gèrent cela automatiquement.
\end{warnbox}

\subsection{Lancer le déploiement}

\begin{enumerate}[label=\arabic*.]
  \item Clique sur \textbf{"Deploy"}
  \item Vercel va compiler et déployer l'application (environ 2 à 4 minutes)
  \item À la fin, tu obtiens une URL publique du type :
        \texttt{https://medicloud-xxxx.vercel.app}
\end{enumerate}

\subsection{Déploiement Continu (CD)}

Après la configuration initiale, chaque fois que quelqu'un push du code
sur la branche \texttt{main} de GitHub, Vercel \textbf{redéploie automatiquement}
l'application sans aucune action manuelle.

\begin{lstlisting}[language=bash]
# Workflow de mise a jour standard
git add .
git commit -m "feat: nouvelle fonctionnalite"
git push origin main
# -> Vercel detecte le push et relance le deploiement automatiquement
\end{lstlisting}

% ─────────────────────────────────────────────────────────────────────────────
\section{Étape 5 — Vérification Post-Déploiement}
% ─────────────────────────────────────────────────────────────────────────────

\subsection{Checklist de vérification}

\begin{itemize}[label=$\checkmark$]
  \item L'URL Vercel s'ouvre et affiche la page d'accueil
  \item La page \texttt{/login} fonctionne et permet la connexion
  \item La recherche de médecins (\texttt{/patient/search}) affiche des résultats
  \item Dans l'onglet \textbf{Cron} du Dashboard Vercel, les tâches planifiées
        sont visibles et actives
  \item Dans Chrome DevTools $\rightarrow$ Application $\rightarrow$ Service Workers, le worker
        est actif (\textit{"Activated and running"})
\end{itemize}

\subsection{Tester les Cron Jobs manuellement}

Pour vérifier que les tâches automatiques fonctionnent, appelle ces URLs
(avec un token d'autorisation) :

\begin{lstlisting}
# Rappels email (simulation)
https://ton-app.vercel.app/api/cron/reminders?bypass=true

# Nettoyage automatique de la BDD
https://ton-app.vercel.app/api/cron/cleanup?bypass=true
\end{lstlisting}

% ─────────────────────────────────────────────────────────────────────────────
\section{Ce que tu Devrais Apprendre — Feuille de Route}
% ─────────────────────────────────────────────────────────────────────────────

Pour comprendre en profondeur comment ce projet fonctionne et progresser
en développement web Cloud, voici les compétences clés à acquérir :

\subsection{Fondations (indispensables)}

\begin{itemize}
  \item \textbf{JavaScript/TypeScript} — Le langage de base de tout le projet.
        Ressource : \href{https://javascript.info}{\texttt{javascript.info}}
  \item \textbf{React.js} — La bibliothèque qui construit l'interface.
        Ressource : \href{https://react.dev}{\texttt{react.dev}}
  \item \textbf{Next.js} — Le framework Full-Stack qui orchestre tout.
        Ressource : \href{https://nextjs.org/learn}{\texttt{nextjs.org/learn}}
  \item \textbf{Git \& GitHub} — Contrôle de version et collaboration.
        Ressource : \href{https://learngitbranching.js.org}{\texttt{learngitbranching.js.org}}
\end{itemize}

\subsection{Back-end \& Bases de Données}

\begin{itemize}
  \item \textbf{SQL \& PostgreSQL} — Comprendre les requêtes, les tables, les jointures.
  \item \textbf{Supabase} — BaaS (Backend as a Service) : Auth, BDD, Storage.
        Ressource : \href{https://supabase.com/docs}{\texttt{supabase.com/docs}}
  \item \textbf{API REST} — Comment concevoir et consommer des APIs HTTP.
  \item \textbf{Authentification JWT} — Tokens, sessions, sécurité.
\end{itemize}

\subsection{Cloud \& DevOps}

\begin{itemize}
  \item \textbf{Vercel} — Déploiement, CI/CD, Serverless Functions, Cron Jobs.
        Ressource : \href{https://vercel.com/docs}{\texttt{vercel.com/docs}}
  \item \textbf{Variables d'environnement} — Gestion des secrets en production.
  \item \textbf{DNS \& Domaines} — Comment configurer un domaine personnalisé.
  \item \textbf{Redis (Upstash)} — Cache distribué pour la scalabilité.
\end{itemize}

\subsection{Concepts Cloud Native avancés (pour aller plus loin)}

\begin{itemize}
  \item \textbf{Progressive Web Apps (PWA)} — Service Workers, Manifestes, Mode hors-ligne.
  \item \textbf{Serverless Architecture} — Functions as a Service (FaaS), Cron Jobs.
  \item \textbf{Edge Computing} — Exécuter du code au plus près de l'utilisateur.
  \item \textbf{Row Level Security (RLS)} — Sécurité au niveau des lignes de BDD.
  \item \textbf{CI/CD Pipelines} — Déploiement automatisé à chaque commit.
\end{itemize}

% ─────────────────────────────────────────────────────────────────────────────
\section{Ressources et Références}
% ─────────────────────────────────────────────────────────────────────────────

\begin{center}
\begin{tabular}{lll}
  \toprule
  \textbf{Ressource} & \textbf{Lien} & \textbf{Usage} \\
  \midrule
  Code source & \href{https://github.com/Noamane2023/medicloud}{\texttt{github.com/Noamane2023/medicloud}} & Accès au projet \\
  Next.js Docs & \href{https://nextjs.org/docs}{\texttt{nextjs.org/docs}} & Framework \\
  Supabase Docs & \href{https://supabase.com/docs}{\texttt{supabase.com/docs}} & Base de données \\
  Vercel Docs & \href{https://vercel.com/docs}{\texttt{vercel.com/docs}} & Déploiement \\
  Upstash Docs & \href{https://docs.upstash.com}{\texttt{docs.upstash.com}} & Redis Cache \\
  Resend Docs & \href{https://resend.com/docs}{\texttt{resend.com/docs}} & Emails \\
  \bottomrule
\end{tabular}
\end{center}

% ─────────────────────────────────────────────────────────────────────────────
\section*{Conclusion}
% ─────────────────────────────────────────────────────────────────────────────

En suivant ce guide, tu as pu :
\begin{itemize}
  \item Installer l'environnement de développement complet
  \item Lancer MediCloud localement sur ton ordinateur
  \item Déployer l'application en production sur Vercel
  \item Comprendre la feuille de route pour maîtriser ces technologies
\end{itemize}

\begin{successbox}{Bravo !}
  Le développement Cloud Native est l'une des compétences les plus demandées
  dans l'industrie tech. Les technologies utilisées dans ce projet (Next.js,
  Supabase, Vercel, Redis) sont utilisées par des milliers d'entreprises
  à travers le monde. Continue à explorer !
\end{successbox}

\vspace{1cm}
\begin{center}
  \textit{Document rédigé avec \LaTeX\ par Noaman Makhlouf — ENSEM 2026}
\end{center}

\end{document}
